Question Four (A)
Effect of Temperature on Charge Carriers in an Extrinsic Semiconductor
When the temperature of an extrinsic semiconductor (n-type or p-type) increases, thermal energy disrupts the crystal lattice, breaking covalent bonds and generating electron-hole pairs.
Minority Charge Carriers: The number of minority carriers is highly sensitive to temperature changes. Because their baseline concentration at room temperature is very low, the thermal generation of new electron-hole pairs causes their concentration to increase exponentially with temperature.
Majority Charge Carriers: The number of majority carriers is nearly unaffected across normal operating temperatures. Because their concentration is primarily determined by the doping level (which is much larger than the thermally generated carriers), the addition of a few thermal carriers causes a negligible percentage change.
(Note: At extremely high temperatures, thermal generation eventually dominates doping, making the semiconductor act as an intrinsic material).
Question Four (B)
i. State whether it is a silicon or germanium diode
This is a silicon diode.
Reasoning: Looking at the data table, the forward current remains near zero ($0\text{ mA}$ and $0.02\text{ mA}$) at voltages up to $0.4\text{ V}$. A significant, sharp conduction begins around $0.6\text{ V}$ ($1.0\text{ mA}$) and rapidly rises by $0.8\text{ V}$ ($100\text{ mA}$). This matches the typical knee voltage (barrier potential) of $0.7\text{ V}$ for silicon, whereas a germanium diode would conduct significantly earlier around $0.3\text{ V}$.
ii. Determine the value of the resistor to limit the forward current to 100mA on a 3 V supply
To design the limiting circuit, we apply Kirchhoff's Voltage Law (KVL) around the series loop:
$$V_{\text{supply}} = V_D + (I_F \times R)$$ 
Identify parameters from the table at $I_F = 100\text{ mA}$:
Target Forward Current, $I_F = 100\text{ mA} = 0.1\text{ A}$
Corresponding Forward Diode Voltage drop from the table, $V_D = 0.8\text{ V}$
Supply Voltage, $V_{\text{supply}} = 3\text{ V}$
Rearrange the formula to find the series resistance ($R$):
$$R = \frac{V_{\text{supply}} - V_D}{I_F}$$ 
$$R = \frac{3\text{ V} - 0.8\text{ V}}{0.1\text{ A}} = \frac{2.2\text{ V}}{0.1\text{ A}} = \mathbf{22\ \Omega}$$ 
